% GENERATED FILE. Edit canonical lessons, then run npm run generate:books.
% canonical-source-hash: fnv1a64:f14c84485c3b07f0
% canonical-lessons: PT-C287-roubo, PT-C287-furto, PT-C287-furtar, PT-C287-carteirista, PT-C287-desaparecimento

\chapter{Roubos}
\label{ch:pt-roubos}
\hlchaptermodality{287}

\begin{chapteropening}
\textbf{By the end of this chapter:} \emph{I can talk about robbery and theft, pickpockets, and things that went missing.}

Complete the last lesson of chapter 287: I can talk about robbery and theft, pickpockets, and things that went missing.
\end{chapteropening}

\section[o roubo]{o roubo — the robbery}

\label{lesson:PT-C287-roubo}



\begin{quote}
Before the new one: say the Portuguese for the victim, then the Portuguese for the homicide.
\end{quote}

\subsection*{You'll want to know: o roubo}
\textbf{o roubo} — \textquotedblleft{}the robbery\textquotedblright{}.

\begin{quote}
Houve um roubo na loja ao lado.
\end{quote}

\begin{quote}
\emph{There was a robbery in the shop next door.}
\end{quote}

Taking something by force or threat.

\begin{cousinweb}
\emph{roubo} is built on \emph{roubar} \textquotedblleft{}to rob\textquotedblright{}, which EtymDB traces through Medieval Latin to Frankish \emph{raubōn} \textquotedblleft{}to plunder\textquotedblright{}.
\end{cousinweb}

\begin{grammarlens}[title={What you've built}]
One more word for telling, explaining and describing at B1.
\end{grammarlens}

\subsection*{Your turn}
\begin{itemize}
\raggedright
  \item \emph{Say it:} \emph{o roubo}
  \item \emph{Say it:} \emph{o roubo}, once more
  \item \emph{Say it:} a sentence of your own with \emph{o roubo}
  \item \emph{Recall:} say the Portuguese for the victim, then the Portuguese for the homicide, then say \emph{o roubo} again
\end{itemize}

\begin{tcolorbox}[breakable,colback=teal!4,colframe=teal!35!black,title={Before you move on}]
What does o roubo mean? (\textquotedblleft{}The robbery\textquotedblright{}.) Say it once more.
\end{tcolorbox}
\section[o furto]{o furto — the theft (without violence)}

\label{lesson:PT-C287-furto}



\begin{quote}
Before the new one: say the Portuguese for the homicide, then the Portuguese for the robbery.
\end{quote}

\subsection*{You'll want to know: o furto}
\textbf{o furto} — \textquotedblleft{}the theft (without violence)\textquotedblright{}.

\begin{quote}
O furto da bicicleta aconteceu de noite.
\end{quote}

\begin{quote}
\emph{The bike was stolen at night.}
\end{quote}

\emph{Roubo} involves force or threat; \emph{furto} is taking something quietly.

\begin{cousinweb}
\emph{furto} is inherited from Latin \emph{fūrtum} \textquotedblleft{}theft\textquotedblright{} (EtymDB), from \emph{fūr} \textquotedblleft{}a thief\textquotedblright{}.
\end{cousinweb}

\begin{grammarlens}[title={What you've built}]
One more word for telling, explaining and describing at B1.
\end{grammarlens}

\subsection*{Your turn}
\begin{itemize}
\raggedright
  \item \emph{Say it:} \emph{o furto}
  \item \emph{Say it:} \emph{o furto}, once more
  \item \emph{Say it:} a sentence of your own with \emph{o furto}
  \item \emph{Recall:} say the Portuguese for the homicide, then the Portuguese for the robbery, then say \emph{o furto} again
\end{itemize}

\begin{tcolorbox}[breakable,colback=teal!4,colframe=teal!35!black,title={Before you move on}]
What does o furto mean? (\textquotedblleft{}The theft (without violence)\textquotedblright{}.) Say it once more.
\end{tcolorbox}
\section[furtar]{furtar — to steal (quietly), to pinch}

\label{lesson:PT-C287-furtar}



\begin{quote}
Before the new one: say the Portuguese for the robbery, then the Portuguese for the theft.
\end{quote}

\subsection*{You'll want to know: furtar}
\textbf{furtar} — \textquotedblleft{}to steal (quietly), to pinch\textquotedblright{}.

\begin{quote}
Alguém furtou a minha carteira no metro.
\end{quote}

\begin{quote}
\emph{Someone pinched my wallet on the underground.}
\end{quote}

Without force; with force, \emph{roubar}.

\begin{cousinweb}
\emph{furtar} is built on \emph{furto} \textquotedblleft{}theft\textquotedblright{} (EtymDB), from Latin \emph{fūrtum}, from \emph{fūr} \textquotedblleft{}a thief\textquotedblright{}.
\end{cousinweb}

\begin{grammarlens}[title={What you've built}]
One more word for telling, explaining and describing at B1.
\end{grammarlens}

\subsection*{Your turn}
\begin{itemize}
\raggedright
  \item \emph{Say it:} \emph{furtar}
  \item \emph{Say it:} \emph{furtar}, once more
  \item \emph{Say it:} a sentence of your own with \emph{furtar}
  \item \emph{Recall:} say the Portuguese for the robbery, then the Portuguese for the theft, then say \emph{furtar} again
\end{itemize}

\begin{tcolorbox}[breakable,colback=teal!4,colframe=teal!35!black,title={Before you move on}]
What does furtar mean? (\textquotedblleft{}To steal (quietly), to pinch\textquotedblright{}.) Say it once more.
\end{tcolorbox}
\section[o carteirista]{o carteirista — the pickpocket}

\label{lesson:PT-C287-carteirista}



\begin{quote}
Before the new one: say the Portuguese for the theft, then the Portuguese for to steal.
\end{quote}

\subsection*{You'll want to know: o carteirista}
\textbf{o carteirista} — \textquotedblleft{}the pickpocket\textquotedblright{}.

\begin{quote}
Um carteirista tirou-me o telemóvel no elétrico.
\end{quote}

\begin{quote}
\emph{A pickpocket took my mobile on the tram.}
\end{quote}

Built on \emph{carteira} ‘wallet’. In Brazil: \emph{o batedor de carteira}.

\begin{cousinweb}
\emph{carteirista} is built on \emph{carteira} \textquotedblleft{}wallet\textquotedblright{}, from \emph{carta}, which comes from Latin \emph{charta} \textquotedblleft{}paper\textquotedblright{} (EtymDB), with \emph{-ista} for the person.
\end{cousinweb}

\begin{grammarlens}[title={What you've built}]
One more word for telling, explaining and describing at B1.
\end{grammarlens}

\subsection*{Your turn}
\begin{itemize}
\raggedright
  \item \emph{Say it:} \emph{o carteirista}
  \item \emph{Say it:} \emph{o carteirista}, once more
  \item \emph{Say it:} a sentence of your own with \emph{o carteirista}
  \item \emph{Recall:} say the Portuguese for the theft, then the Portuguese for to steal, then say \emph{o carteirista} again
\end{itemize}

\begin{tcolorbox}[breakable,colback=teal!4,colframe=teal!35!black,title={Before you move on}]
What does o carteirista mean? (\textquotedblleft{}The pickpocket\textquotedblright{}.) Say it once more.
\end{tcolorbox}
\section[o desaparecimento]{o desaparecimento — the disappearance}

\label{lesson:PT-C287-desaparecimento}



\begin{quote}
Before the new one: say the Portuguese for to steal, then the Portuguese for the pickpocket.
\end{quote}

\subsection*{You'll want to know: o desaparecimento}
\textbf{o desaparecimento} — \textquotedblleft{}the disappearance\textquotedblright{}.

\begin{quote}
O desaparecimento do quadro foi notícia em todo o mundo.
\end{quote}

\begin{quote}
\emph{The painting’s disappearance was news all over the world.}
\end{quote}

Built on \emph{desaparecer}.

\begin{cousinweb}
\emph{desaparecimento} is built on \emph{desaparecer} \textquotedblleft{}to disappear\textquotedblright{}, from \emph{des-} and \emph{aparecer}, which goes back to Latin \emph{appārēre} \textquotedblleft{}to come into sight\textquotedblright{} (Lewis \& Short).
\end{cousinweb}

\begin{grammarlens}[title={What you've built}]
One more word for telling, explaining and describing at B1.
\end{grammarlens}

\subsection*{Your turn}
\begin{itemize}
\raggedright
  \item \emph{Say it:} \emph{o desaparecimento}
  \item \emph{Say it:} \emph{o desaparecimento}, once more
  \item \emph{Say it:} a sentence of your own with \emph{o desaparecimento}
  \item \emph{Recall:} say the Portuguese for to steal, then the Portuguese for the pickpocket, then say \emph{o desaparecimento} again
\end{itemize}

\begin{tcolorbox}[breakable,colback=teal!4,colframe=teal!35!black,title={Before you move on}]
What does o desaparecimento mean? (\textquotedblleft{}The disappearance\textquotedblright{}.) Say it once more.
\end{tcolorbox}
